\section{Data and assumptions}
\label{sec:data}

Each Set~A instance is given by three JSON files.
The network file provides vertices, directed links, administrative metrics $\omega_a$ and
capacities $c_a$.
The traffic file provides, for every demand~$d$, a source $s(d)$, a target $t(d)$ and a
volume $v_{dt}$ in each time slot.
The scenario file gives $\mathrm{maxSeg}=6$, the period-1 budget $\kappa$ and the links
removed at period~1.
Period~0 has no intervention and no budget.
Identifiers are stored as given; they are not assumed to form an interval, and no Set~A
instance contains two links with the same ordered pair of endpoints.

For a period~$t$, let $G(t)=(V,A(t))$ be the directed graph of available links.
Dijkstra's algorithm from every vertex returns the distance $d_{ij}(t)$ and the number
$\sigma_{ij}(t)$ of shortest directed $i$--$j$ paths (path counts are stored in logarithmic
form).
For an arc $a=(u,v)$, the course note~\citep{courser2026} defines the fraction of those
paths that use~$a$:
\begin{equation}
r_{ijat}=
\begin{cases}
\dfrac{\sigma_{iu}(t)\,\sigma_{vj}(t)}{\sigma_{ij}(t)},
  & \text{if } d_{iu}(t)+\omega_a+d_{vj}(t)=d_{ij}(t),\\[0.35em]
0, & \text{otherwise.}
\end{cases}
\label{eq:r}
\end{equation}
This uniform-over-path formula is the one used throughout the experiments.
It is not identical to a router-local ECMP recursion, which can assign different probabilities
to complete shortest paths on an asymmetric DAG; loads are therefore comparable within this
report, not necessarily to an official challenge evaluator that uses the other rule.
If $j$ is unreachable from $i$, the segment $(i,j)$ is omitted.
